\section{Introduction}

The cerebral cortex performs perception, action, and internal-model inference by coordinating activity across distributed areas. This long-range coordination is thought to be mediated by specific classes of pyramidal neurons whose axonal projections span the cortical mantle, and its disruption has been repeatedly invoked as a substrate of psychotic disorders~\cite{uhlhaas2015oscillations, uhlhaas2013highfrequency, bahner2017hippocampal}. Hierarchical, anatomically graded connectivity between cortical areas provides the scaffold on which such coordination is built~\cite{harris2019hierarchical, siegle2021survey}, and different classes of glutamatergic neurons contribute distinctly to inter-areal and cortico-cortical signaling~\cite{matho2021genetic, biccn2021multimodal, liu2024celltype}. Despite decades of progress on the receptor pharmacology of antipsychotic drugs, how their therapeutic action reshapes this distributed cortical circuit remains poorly understood.

The dominant pharmacological account of antipsychotic efficacy emphasizes dopamine D2 receptor antagonism, with complementary actions at 5-HT2A and other monoaminergic targets~\cite{howes2009dopamine, kambeitz2014alterations}. While these frameworks have guided drug development, they do not specify where in the cortex antipsychotics change neural activity, nor whether their effect is uniform across cortical layers and cell types. Because positive symptoms of psychosis are associated with aberrant salience and disordered long-range synchrony~\cite{uhlhaas2015oscillations, bahner2017hippocampal}, a cortical circuit account of antipsychotic action must go beyond receptor occupancy and ask how receptor binding translates into changes in inter-areal coupling.

Layer 5 intratelencephalic (L5 IT) pyramidal neurons are candidate substrates for such coupling. They project densely within cortex and between hemispheres, and recent work using genetic labeling has delineated their anatomical reach and functional contributions to cortico-cortical signaling~\cite{liu2024celltype, shinotsuka2023layer5, atudorei2024bilateral, matho2021genetic}. The Tlx3-Cre driver line selectively targets L5 IT neurons, and in combination with reporter lines such as Ai148 that express GCaMP6 under Cre control~\cite{daigle2018suite, chen2013ultrasensitive}, it enables cell-type- and layer-specific readouts of distributed cortical activity. Complementary Cre lines such as Cux2-CreERT2, combined with tamoxifen-induced recombination, enable comparable access to layer 2/3 pyramidal neurons~\cite{daigle2018suite}. Pan-neuronal expression can be achieved via the retro-orbitally administered AAV-PHP.eB capsid~\cite{chan2017engineered}, providing a third comparison population.

Widefield calcium imaging across the dorsal cortex of head-fixed mice has emerged as a powerful approach for capturing distributed activity at a scale commensurate with inter-areal coupling~\cite{musall2019single, leinweber2017sensorimotor}. Combined with head-fixed virtual reality, it permits the alignment of inter-areal activity to behaviorally well-defined events such as locomotion onsets and visual-flow onsets~\cite{leinweber2017sensorimotor, keller2012sensorimotor, attinger2017visuomotor, keller2018predictive}. Under a predictive-processing framing, these events separate the bottom-up and top-down components of cortical activity and thus provide a behaviorally anchored readout through which drug-induced changes in coupling can be interpreted~\cite{keller2018predictive, zmarz2016mismatch}. Methodological refinements---including hierarchical bootstrap statistics for nested neuronal data~\cite{saravanan2020hierarchical} and percentile-based baseline correction adapted from two-photon imaging~\cite{dombeck2007imaging}---allow this approach to be combined with pharmacological perturbations within individual mice.

In this work, we ask whether antipsychotic drugs selectively alter long-range cortical coupling, and whether this alteration is specific to a molecularly defined subclass of pyramidal neurons. We combine widefield calcium imaging across twelve regions of dorsal cortex with cell-type-specific expression of GCaMP6 in pan-neuronal (C57BL/6 with AAV-PHP.eB), L5 IT (Tlx3-Cre~$\times$~Ai148), and L2/3 (Cux2-CreERT2~$\times$~Ai148) populations, before and after single intraperitoneal injections of clozapine, aripiprazole, haloperidol, amphetamine, or saline. We report three main findings: (i) antipsychotics reduce short- and long-range correlations in L5 IT neurons substantially more than in pan-neuronal recordings; (ii) the effect is substantially weaker in L2/3 Cux2 neurons than in L5 IT Tlx3 neurons; and (iii) the psychostimulant amphetamine, which shares dopaminergic but not antipsychotic properties, does not reproduce the decorrelation pattern. Together, these observations identify L5 IT neurons as a cell-type-specific cortical substrate of antipsychotic action and support a circuit-level account of therapeutic efficacy.

\section{Background and Related Work}

\paragraph{Cortical interactions and psychosis.}
The idea that psychosis involves altered long-range cortical coordination is supported by both electrophysiological and hemodynamic evidence. Abnormal beta and gamma oscillations, as well as reduced large-scale synchrony, have been reported repeatedly in schizophrenia~\cite{uhlhaas2015oscillations, uhlhaas2013highfrequency}. Hippocampal--prefrontal coupling during cognitive tasks has been proposed as a translational biomarker whose alterations track illness and genetic risk~\cite{bahner2017hippocampal}. These observations motivate measurements of inter-areal correlation in animal models of pharmacological intervention.

\paragraph{Dopaminergic and receptor-level accounts.}
The dopamine hypothesis of schizophrenia, and the subsequent aberrant salience framework, anchors antipsychotic action in D2 receptor antagonism and associated dopamine dysfunction~\cite{howes2009dopamine, kambeitz2014alterations}. Drugs such as clozapine, aripiprazole, and haloperidol differ in their affinities for D2, 5-HT2A, and other targets, yet all are clinically effective against positive symptoms. This clinical convergence despite pharmacological divergence suggests a shared downstream effect at the circuit level that receptor-centric models do not fully specify.

\paragraph{Genetic and anatomical dissection of cortical cell types.}
Work over the past two decades has produced Cre driver lines and anatomical atlases that enable cell-type-specific access to cortical pyramidal neuron subclasses. L5 IT neurons, labeled by Tlx3-Cre, form dense cortico-cortical and cortico-striatal projections~\cite{liu2024celltype, matho2021genetic, biccn2021multimodal, shinotsuka2023layer5, atudorei2024bilateral}, and their activity is implicated in coordinated behaviors that require distributed cortical signals. Layer 2/3 neurons in sensory cortex, by contrast, participate in local computations and are targeted by the Cux2-CreERT2 line~\cite{daigle2018suite}. Hierarchically organized connectivity across cortical areas and between cortex and thalamus provides the structural backbone for the interactions measured here~\cite{harris2019hierarchical, siegle2021survey}.

\paragraph{Widefield imaging of cortical dynamics.}
Widefield calcium imaging of GCaMP indicators across dorsal cortex captures distributed activity at millisecond-to-second timescales and has been used to relate inter-areal signals to behavior~\cite{musall2019single, leinweber2017sensorimotor}. Genetically encoded calcium sensors, in particular GCaMP6 variants~\cite{chen2013ultrasensitive}, and Cre-dependent transgenic reporter lines such as Ai148~\cite{daigle2018suite, bounds2023alloptical}, make cell-type-specific widefield recordings practical. Retro-orbital delivery of AAV-PHP.eB enables pan-neuronal or Cre-dependent expression without intracranial injection~\cite{chan2017engineered}. Together, these tools allow the same imaging pipeline to be applied across genetically defined populations, providing a controlled comparison of drug effects across cell types.

\paragraph{Predictive processing and sensorimotor integration.}
A complementary line of work has framed cortical activity during locomotion and visual flow in terms of predictive processing~\cite{keller2012sensorimotor, keller2018predictive, zmarz2016mismatch, attinger2017visuomotor, leinweber2017sensorimotor}. Closed-loop and open-loop locomotion onsets, as well as mismatches between predicted and actual visual flow, reveal separable bottom-up and top-down contributions to cortical activity. In particular, top-down projections from frontal and cingulate areas shape activity in primary visual cortex and carry signals consistent with predictions of visual flow based on motor output~\cite{leinweber2017sensorimotor}. These behaviorally defined events serve here as internal benchmarks for comparing naive and drug-treated activity.

\paragraph{Statistical analysis of nested neural data.}
Widefield correlation analyses yield nested data structures in which pairs of regions are grouped within mice, and mice within genotypes. Traditional tests under-report variance in such data. The hierarchical bootstrap approach~\cite{saravanan2020hierarchical} provides confidence intervals that honor this nesting and is used throughout for time-course estimates. For discrete correlation comparisons, non-parametric rank-sum tests with family-wise error correction are used, as described in Methods. Baseline correction and slow-drift removal follow established two-photon practice adapted to widefield data~\cite{dombeck2007imaging}.

In summary, prior work provides the tools (cell-type-specific GCaMP6 expression, widefield imaging, robust statistics) and the conceptual frame (long-range cortical coupling, predictive processing, dopaminergic pharmacology) within which the present experiments are situated. What has been missing is a direct, cell-type-resolved, within-mouse comparison of inter-areal correlations before and after antipsychotic administration. The experiments reported below provide such a comparison and identify L5 IT neurons as the population most strongly affected by antipsychotic treatment.
